\scanpage{81}

% L01 -- top, continuation
i.e. sss on a simultanément (posant $i=n-d$)

% L02 -- top, brace and qualifications
\[
\left\{\begin{aligned}
\mu_i\ (=\lambda_i-1)&=0&&
\begin{gathered}
\text{(i.e. }2\cdot1_k=0\text{ dans le cas des réflexions\unsure{)|,}}\\
\text{i.e. } k\text{ de car. }2)
\end{gathered}\\
\beta_i=\beta_{i-1}&=0&&
\begin{gathered}
\text{(se réduit à la condition }\beta_i=0\\
\text{si }i=0,\text{ et à la condition}\\
\beta_{i-1}=0\text{ si }i=n)
\end{gathered}
\end{aligned}\right.
\]

% L03 -- upper left, abandoned headings
\deleted{\uline{Cor} \unsure{D'autre} \uline{Dém.}}

% L04 -- upper middle, corollary a
\uline{Cor.} \textcircled{a} Il y a au plus $2(n+2)$ sous-espaces projectifs invariants par $G$.

% L05 -- upper middle, b
\textcircled{b} Pour qu'il y en ait $2(n+2)$, il f et s. que $\beta_i=0$ $(0\le i\le n-1)$ et $\mu_i\ne0$ \deleted{\unsure{$\forall$}} $(0\le i\le n)$ (\unsure{cette condition|cette autre} signifie aussi, dans le cas des réflexions, que car. $k\ne2$).

% L06 -- middle, c
\textcircled{c} \deleted{Pour} Il y a au moins les quatre ss-espaces invariants $\varnothing,X,\{c\},C$. Pour que ce soient les seuls, il \deleted{\unsure{f et s.|faut et}} \inserted{suffit} que $\beta_i\ne0\qquad 0\le i\le n-1$ (ce qui implique que les $2(n+\unsure{1|2})$ espaces $Z_i,Z'_i$ sont distincts). Il faut \inserted{et s.} que $\beta_i\ne0\quad 0\unsure{\le|<}i<n-1$, et que $(\beta_0\deleted{\unsure{\ne\beta_n|=\beta_n}}=0\Rightarrow\mu_0=0)$, $(\beta_{n-1}=0\Rightarrow\mu_n=0)$ [donc, dans le cas des réflexions, que

% L07 -- lower middle, cases
\[
\left\{\begin{aligned}
&\text{si car. }k\ne2:\ \text{les }\beta_i\ne0\ \inserted{0\le i\le n-1}\quad
(\text{i.e. }\beta=\prod_i\beta_i\ne0\\
&\text{si car }k=2\qquad\text{les }\beta_i\ne0\text{ si }0<i<n-1.
\end{aligned}\right.
\]

% L08 -- lower part, struck d
\deleted{d) Pour que \deleted{\unsure{la sous-représentation|sous-espace}} \inserted{la représentation projective} de $G$ \uline{dans $C$} soit irréductible, il f et s. que $\rho\ne0$ et les $\beta_i\ne\unsure{2|0}$, i.e. que}

% L09 -- lower left, enclosed and linked to d
\margin{Si $c\notin C$ i.e $\rho\ne0$}

% L10 -- foot, new d, continues
d) les sous-espaces projectifs stables de $C$ \deleted{\unsure{et de $\{c\}$ (dans le cas $c\in C$)}} $\ne C$ \deleted{et}, de $\varnothing$ sont en corr. 1-1 avec les $0\le i\le n-1$ tels que $\beta_i=0$, \deleted{et} un tel $i$ correspondant $Z_i^{\deleted{\unsure{\prime|*}}}$ (de dim $n-1-i$). Ces sous
% archive sheet number: 80
